% Audited against 20260905_100500_required_target_depth_diagnostic.
% This interpretation must be reviewed if the diagnostic data are regenerated.
All six depths complete the same nine of 12 validation tasks; the other three lack visibility of at least one required item. On this common nine-task cohort, increasing depth from two to three reduces mean graph cost by 11.19\%, from 133.450 to 118.520. Depth three captures 84.11\% of the observed reduction from depth one to the lowest mean cost among depths one through six. Depth four further reduces mean cost by 2.74\%, to 115.272, but raises median planning time from \unit{1.009}{s} to \unit{4.676}{s} and mean safety-screen count by a factor of 4.74. Depths five and six yield no further cost or completion improvement, while their median times increase to \unit{17.660}{s} and \unit{53.013}{s}, respectively. These are 17.50 and 52.52 times the depth-three median; the corresponding screen-count ratios are 17.86 and 54.15.

These observations support retaining the frozen depth-three policy as a computational compromise, not as the minimum-cost or uniquely optimal depth. If the additional computation is acceptable, depth four provides lower cost in this diagnostic. The plateau at depths four through six is empirical, not an exact-optimality certificate, and realized receding-policy costs need not decrease monotonically with depth.
